Operating the harness of \S\ref{sec:harness} on real hardware turned the
design rationale of \S\ref{sec:harness-rationale} from an argument into
an observation: six recurring geometry-grounding bugs, each traced to
its root cause using the harness's own per-run debug artifacts (raw
model response, config, and, in VLM mode, the captured camera frame).
\rev{These bugs surfaced on the same physical rig and task set used
throughout this paper: a UR5cb manipulator with an OnRobot RG2 gripper
and a calibrated RGB-D camera, operating over a real scene of four
color-distinguished cubes and a tray, executing pick-and-place, sorting,
and arithmetic value-mapping tasks (\S\ref{sec:evaluation} details the
full protocol).} Table~\ref{tab:taxonomy} summarizes all six. Below, we
develop the two most instructive of the five closed cases in full, since
both recurred independently in two different planning pipelines and so
best show that the pattern is structural, not a one-off slip, then treat
the sixth, still-open bug on its own (\S\ref{sec:case-collision}), which
turns out to be the dominant failure mode in \S\ref{sec:evaluation}.

\begin{table*}[t]
\centering
\caption{Geometry-grounding failure taxonomy. ``Sensor/actuator required''
marks the physical property that makes each bug impossible to observe in
a symbolic or scripted simulator.}
\label{tab:taxonomy}
\begin{tabularx}{\textwidth}{@{}p{3.1cm}Xp{3.9cm}@{}}
\toprule
\textbf{Bug} & \textbf{Symptom and root cause} & \textbf{Sensor/actuator required} \\
\midrule
Top-vs-mid-body grasp/release mismatch &
  Release-lift height set to an object's full height, implicitly assuming
  a top-surface grasp; the executor actually descends mid-body
  ($\descent = \min(0.5\height, \descent_{\max})$). Net effect: a
  systematic release overshoot of roughly half the object's height.
  Recurred independently in two pipelines (VLM mode, then LLM mode). &
  Real gripper kinematics (mid-body descent is a deliberate stability
  choice, not a simulator default). \\
\addlinespace
Double-counted stacking height &
  In a VLM$\to$LLM hybrid pipeline, a depth-measured true top surface had
  the VLM's separate, non-depth-verified size estimate added on top of it
  a second time by the downstream LLM's own stacking formula
  ($\text{position}.z + \text{size}[2]$). &
  Real depth measurement of an already-correct surface height. \\
\addlinespace
Footprint- vs.\ point-distance target matching &
  A release genuinely inside a tray's real bounding box, but far from its
  registered center point, was treated as having missed the target zone
  by a point-distance check. &
  Real target geometry with non-trivial spatial extent (tens of
  centimeters), not a point. \\
\addlinespace
Grasp-width sampled from the silhouette edge &
  Grasp width derived by deprojecting a bounding box's left/right edge
  pixels, exactly on the object's silhouette boundary, the single
  worst location for valid depth (occlusion/parallax). One failed edge
  sample aborted an otherwise valid plan. &
  Real depth noise at object boundaries; does not exist for a rendered or
  symbolic bounding box. \\
\addlinespace
Width-dependent TCP offset &
  The gripper's fingers pivot, so flange-to-contact distance is a
  function of how far apart they are commanded to open, not a constant,
  so it is wrong for any object whose width diverges from a single
  calibrated offset. &
  Real gripper mechanism (finger kinematics), absent from a symbolic
  action space. \\
\addlinespace
Multi-object sequential-placement collision \textbf{(open)} &
  Zone-aware spacing keeps a target zone's \emph{final} release positions
  from overlapping, but does not prevent the arm clipping an
  already-placed object en route to the next release, nor a released
  object settling into a neighbor from residual momentum or a few cm of
  height error. Dominant failure mode in a 120-trial re-evaluation: every
  non-empty operator note (22/22) describes a collision or a
  fallen/displaced object during a multi-object task, in \textbf{both}
  the LLM and VLM pipelines. &
  Real multi-body contact between two already-placed objects and the
  arm's approach trajectory, absent from any single-object placement
  check, symbolic or real. \\
\bottomrule
\end{tabularx}
\end{table*}

\subsection{Case 1: top-vs-mid-body grasp/release mismatch}
\label{sec:case-grasp-mismatch}

Our executor deliberately grasps objects mid-body rather than at their
top surface: the pick descent is $\descent = \min(0.5\height,
\descent_{\max})$, where $\height$ is the object's height and
$\descent_{\max}$ a clamp preventing the gripper's rigid mounting body
from crashing into tall objects, since centering the grasp mid-body
produces a more stable hold than pinching at the top edge.

The paired \emph{release} action, however, computes how far to lift the
object before opening the gripper (\code{object\_height}, added to
\code{release\_position.z} at execution time) by assuming it was grasped
at the top, wrong by construction whenever the pick descended mid-body,
which is always. The result is a systematic overshoot of roughly
$0.5\height$: objects land, or are dropped from, a height proportional
to half their own size rather than fractions of a centimeter.

We fixed this once, in the VLM pipeline, by deriving
\code{object\_height} from measured depth
($\height = \max(\surfz - \topz, \height_{\min})$) instead of a
model-reported size. It then recurred independently in the LLM pipeline
weeks later, because that pipeline's prompt still asked the model to
compute \code{object\_height} as \code{size[2]}, the same assumption the
VLM-mode fix had already identified as wrong. Fixing the bug once, in
one pipeline, did not fix it everywhere it existed; the general fix
(\S\ref{sec:grounding-layer}) computes this quantity deterministically
for \emph{every} pipeline, since the error is not a mistake any
particular prompt makes but a mismatch between what the model is asked
to assume and what the gripper actually does.

\subsection{Case 2: double-counted stacking height}
\label{sec:case-double-count}

Our static scene description uses the convention
\code{position.z} $=$ base reference, \code{size[2]} $=$ height, so
\code{position.z} $+$ \code{size[2]} reconstructs a target's top
surface, the formula every LLM-mode prompt uses for a stacking release
height. In our VLM$\to$LLM hybrid pipeline, a VLM grounding call first
deprojects a target's pixel location to a depth-measured world
coordinate, then hands the resulting scene to the same LLM planner used
in static-scene mode. An early version of this hybrid set the generated
scene's \code{target.position.z} directly to that depth-measured top
surface, already correct, while leaving the VLM's separate,
un-verified \code{size[2]} guess untouched. The downstream planner then
applied its normal stacking formula, adding the guessed height on top of
an already-correct value and inflating every stack-on-target release by
however wrong that guess happened to be.

The fix follows Case 1's pattern: derive the target's real height from
depth ($\height = \max(\topz - \surfz, 0)$, no minimum floor, since a
flat target zone marked on the table should read as near-zero height),
then reset \code{position.z} to a base reference so the existing
downstream formula reconstructs the true top surface instead of
double-counting it. The sensor-computed quantity is inserted
\emph{into} the representation the model already expects, rather than
trusted from the model's own output.

\subsection{Case 6 (open): multi-object sequential-placement collision}
\label{sec:case-collision}

Unlike Cases 1--2, this bug is not closed by \S\ref{sec:grounding-layer}.
When a task places several objects into the same target zone in
sequence, each release's \emph{final} position is spaced to avoid
overlapping the others already placed, by a grid clamped to the zone's
real bounds (\S\ref{sec:grounding-layer}). This closes the
\emph{static} half: no two final positions occupy the same footprint. It
does not address the \emph{dynamic} half: nothing constrains the arm's
trajectory while approaching or releasing the $n$-th object against the
$(n-1)$-th already sitting in the zone, and nothing constrains a
released object's small residual motion from displacing a neighbor
afterward. Both are genuinely three-dimensional, time-dependent contact
problems, not resolvable by correcting a single scalar the way Cases
1--2 were.

We identified this bug from the same debug artifacts as Cases 1--2, but
by the time it was isolated it was already the dominant failure mode of
\S\ref{sec:evaluation} (22 of 22 non-empty operator notes, both
pipelines), rather than something we could close before reporting
results. We report it as an identified, partially-mitigated
\textbf{open} problem, not a claimed contribution;
\S\ref{sec:discussion} returns to what closing it would require.

\subsection{The common pattern}

The six bugs in Table~\ref{tab:taxonomy} share a structure: a quantity
trivially, exactly correct in a symbolic or scripted simulator becomes,
on real hardware, (a) a sensor measurement subject to depth noise and
occlusion, (b) a function of real actuator kinematics a text-based
action schema has no natural way to express, or (c), for Case 6, a
consequence of real multi-body contact over time. None of these three
can occur where object geometry is authored by the experiment designer
or perfectly known by the physics engine, precisely the setting where
LLM/VLM robot planners are most commonly
benchmarked~\cite{favaliRAS}. \rev{\S\ref{sec:grounding-layer} turns this
shared structure into a method: wherever a quantity falls into case (a)
or (b), compute it deterministically instead of trusting the model.}
